% Global extension, independently mathematically audited 2026-10-28.
% Exact polyhedral classification is computer-assisted; see retained verifier.
\section{The sharp global six-simplex inequality}
\label{sec:sixdim-global}

The local example from Theorem~\ref{thm:sixdim-golden-local} is in fact
optimal among all six-simplices. The proof combines a small hypergraph
argument with an exact finite description of a linear-ordering relaxation.
The finite description is independently reproduced with two exact polyhedral
backends; the inequalities used after that description are explicit.

\begin{theorem}\label{thm:sixdim-global}
Every six-dimensional simplex $K$ satisfies
\[
 \frac{c_{\mathrm{EHZ}}(K)^3}{6\operatorname{vol}_6K}
 \le \frac{945+75\sqrt5}{1922}.
\]
Equality holds precisely for the simplex in
Theorem~\ref{thm:sixdim-golden-local}, up to translation, positive homothety,
and linear symplectomorphism.
\end{theorem}

Put $q_*=(63+5\sqrt5)/1922$. By
\eqref{eq:sixdim-simplex-ratio}, it suffices to maximize
$|\operatorname{Pf}\Omega|$ on $C_6$. Let
\[
 D_6=\operatorname{conv}\bigl(\{\pm e_{ij}:1\le i<j\le6\}
       \cup\{\pm t_{ijk}:1\le i<j<k\le6\}\bigr),
\]
so $D_6$ has
seventy rank-two generators and is contained in $C_6$.

\subsection{The rank-two hull}

Introduce seven points $u_1,\ldots,u_6,u_7$ with $u_i=e_i$ for $i\le6$
and $u_7=0$. For each three-element set $T=\{i,j,k\}$, choose the form
\[
 b_T=(u_i-u_k)\wedge(u_j-u_k).
\]
Then $D_6=\operatorname{conv}\{\pm b_T\}$. The orientation choices will
not matter for an upper bound. If two triples share two labels, their forms
share a wedge factor, so $b_T\wedge b_U=0$. Also $b_T\wedge b_T=0$.
Consequently the Pfaffian cubic in coefficients of these forms is
multiaffine, and contains no monomial involving both $T$ and $U$ when
$|T\cap U|\ge2$.

Cancel antipodal coefficients in a convex decomposition. It becomes
$\Omega=\sum_Td_Tb_T$ with $\sum_T|d_T|\le1$. Homogeneity permits a
nonzero global maximizer to be chosen with total mass one. Among global
maximizers choose one with smallest support. If two supported types share
two labels, fix their signs and their total absolute mass, and vary the
split between them. The Pfaffian is affine in that split. Its absolute
value therefore attains a maximum at an endpoint, removing one supported
type without lowering the objective. This contradicts minimality. Thus
some global maximum is supported on a family of triples whose pairwise
intersections have size at most one.

We record the elementary classification needed here. A maximal such family
on seven labels has either five or seven triples. Indeed, each triple uses
three of the twenty-one label pairs, so there are at most seven triples.
No label can have degree zero: otherwise maximality requires all pairs on
the other six labels to be covered, which would give degree $5/2$ at each
of those labels. If a label has degree one, with its only triple $vab$,
all six pairs among the remaining four labels $R$ must be covered.
Without a triple contained in $R$, this requires six additional triples,
forcing all twenty-one pairs to be covered and contradicting degree one.
With a triple $xyz\subset R$, the three pairs $wx,wy,wz$, where
$R=\{w,x,y,z\}$, require three separate triples using $a$ or $b$.
One of the pairs $wa,wb$ would then repeat. Hence every label has degree
at least two, and there are at least five triples. Six triples would leave
three uncovered pairs with even degree at every label, hence an uncovered
triangle, contradicting maximality.

With five triples the degree sequence is $(3,2,2,2,2,2,2)$. The three
triples through the degree-three label match the other six labels. The two
remaining triples partition those six labels, and each matching pair crosses
that partition. Up to relabeling this is
\[
 \{135,246,127,347,567\}.
\]
With seven triples every label pair is covered exactly once: the family is
the Fano plane.

For three members of a linear triple family, the Pfaffian coefficient is
zero if their union has at most six labels, since the forms lie in a
difference space of dimension at most five. When their union has seven
labels, their six edge factors form a spanning tree and the coefficient is
$\pm1$. For the five-triple configuration, taking absolute values therefore
gives the polynomial
\[
 xy(z+u+v)+zuv.
\]
Its maximum under nonnegative total mass one is $q_*$, as computed in
\eqref{eq:sixdim-F32}. For the Fano configuration the unsigned polynomial
is the Fano-plane Lagrangian, because the nonzero triples are precisely
three concurrent Fano lines.

Here is a short bound for the latter polynomial $P$. On the boundary,
remove a zero-weight point. The four remaining Fano triples are transversal
to a partition of the six remaining points into three pairs. Hence $P$ is
at most the product of the three class masses, and is at most $1/27$.
At a full-support maximum, the seven partial derivatives all equal $3P$.
Their sum counts every pair exactly once, giving
\[
 21P=\sum_{i<j}x_ix_j\le\frac37,
 \qquad P\le\frac1{49}.
\]
Thus $P\le1/27<q_*$. Together with the explicit attaining example, this
proves
\begin{equation}\label{eq:sixdim-ranktwo-bound}
 \max_{\Omega\in D_6}|\operatorname{Pf}\Omega|=q_*.
\end{equation}

\subsection{The finite polyhedral decomposition}

All pairings in this subsection use the fifteen upper-triangular entries.
Define
\[
 F=\frac15\begin{pmatrix}0&M\\-M^T&0\end{pmatrix},
 \quad M=2I-\mathbf1\mathbf1^T,
 \qquad
 p_0=\begin{pmatrix}0&N\\-N^T&0\end{pmatrix},
 \quad N=I-\mathbf1\mathbf1^T.
\]
The form $F$ is a normalized three-fence. Let $\mathcal F$ be its orbit
under the seven-label permutation action. The action on forms comes from
relabeling the seven barycentric gradients and has determinant $\pm1$;
it preserves $C_6$, $D_6$ and the absolute Pfaffian. Its dual action on
functionals is contragredient.

The retained replay in
\path{certificates/sixdim-global/README.md} specifies four small programs:
\begin{center}
\begin{tabular}{ll}
\path{algebra.py} & \path{local_maximum.py}\\
\path{full_classification.py} & \path{dual_coverage.py}
\end{tabular}
\end{center}
The latter two use exact cdd and PPL polyhedral backends. Their generated
outputs and a fresh copied-package run receipt are retained beside the
programs. This is a computer-assisted completeness check, not a formalized
proof in a proof assistant.

The exact finite calculations establish the following facts:
\begin{enumerate}
\item $Q_6=D_6^\circ$ has exactly $1560$ vertices: the $720$ order
functionals and the $840$ elements of the dual orbit of $p_0$.
\item $\mathcal F$ has $840$ elements. The canonical pairing is
$p_0\cdot F=6/5$, whereas $p_0\cdot F'\le4/5$ for every other
$F'\in\mathcal F$. Every fence satisfies every order inequality.
\item The vertex $p_0$ has eighteen neighbors in $Q_6$, all order
functionals. They are exactly the eighteen orders with $F\cdot p=1$.
The face $B=\{X\in D_6:p_0\cdot X=1\}$ has eighteen rank-two vertices
and affine dimension fourteen.
\end{enumerate}

For completeness, the completeness check in the first item uses graph
connectivity, not a sample of vertices. Six adjacent-transposition generators
preserve the seventy inequalities of $Q_6$ and generate the two disjoint
orbits. At an order representative, both exact polyhedral backends give
$186$ extreme tangent rays, whose exact endpoints are $165$ order vertices
and $21$ fractional vertices. At $p_0$ they give eighteen rays, all with
order endpoints. Both representatives have active rank fifteen. The candidate
set is therefore closed under every neighbor in the one-skeleton. Since a
bounded polytope has connected one-skeleton, there are no further vertices.

Cut $Q_6$ by the $840$ inequalities $F'\cdot p\le1$. Each cut removes
its corresponding fractional vertex and meets all its incident edges at
existing integral neighbors. This can be done sequentially without creating
vertices: all other fractional vertices lie strictly inside the current cut,
so their tangent cones are unchanged; their original edges to integral
neighbors remain feasible in full. Thus each later cut has the same crossing
edges and introduces no new endpoint. The resulting polytope is precisely
the convex hull of the order functionals. Polarity gives
\[
 C_6=\operatorname{conv}(D_6\cup\mathcal F).
\]

Moreover, the midpoint of any two distinct fences lies in $D_6$. On an
order functional both endpoint scores are at most one; on a fractional
functional at most one is $6/5$, while the other is at most $4/5$.
Replacing equal masses of two fences by their midpoint repeatedly reduces
any convex combination to a point of $D_6$ or a combination of $D_6$ and
one fence. Each fence violates exactly one facet inequality of $D_6$.
Consequently $C_6$ is covered by $D_6$ and the $840$ pyramids
$\operatorname{conv}(F',B_{F'})$. It remains to bound one pyramid.

\subsection{A quadratic bound on a fence pyramid}

Write $X=\sum_{i=1}^{18}x_iv_i\in B$, with $x_i\ge0$ and
$\sum_i x_i=1$. Exact Pfaffian expansion gives
\[
 \operatorname{Pf}(tF+sX)
 =\frac4{125}t^3+\frac2{25}t^2s+Q(X)ts^2
       +\operatorname{Pf}(X)s^3.
\]
Here $Q(X)=x^THx/10$, where $H$ is a symmetric integer matrix with zero
diagonal and entries $0,1,2$. The graph of entries equal to two is the
disjoint union of six three-vertex paths. In each path, the two leaves have
entry zero between them; entries between distinct paths are zero or one.
These explicit coefficient properties give
\[
 \mathbf1\mathbf1^T-H
 =\sum_{\text{six paths }(c,l,r)}
       (e_c-e_l-e_r)(e_c-e_l-e_r)^T+R,
 \qquad R\ge0\ \text{entrywise}.
\]
Hence $0\le Q(X)\le1/10$. By~\eqref{eq:sixdim-ranktwo-bound}, for
$t,s\ge0$,
\[
 |\operatorname{Pf}(tF+sX)|
 \le\frac4{125}t^3+\frac2{25}t^2s+\frac1{10}ts^2+q_*s^3
 \le q_*(t+s)^3.
\]
The last inequality is strict when $t>0$, because $q_*>1/30$.
This proves the global bound, with strictness outside $D_6$.

\subsection{Equality}

Begin with any global maximizer in $D_6$ and apply the signed mass-merging
procedure. At each interior split, the absolute value of an affine function
already has its global maximum; that affine function must be constant.
Thus the procedure gives a finite path of global maximizers to a supported
linear triple packing. The Fano case is too small. In the five-triple case,
equality forces the golden weights and common signs for the two central
variables and for the three outer variables. All four resulting sign
patterns are relabelings of $\Omega_*$: $(1\,3)(2\,4)$ flips the two
triangle signs, and $(1\,2)(3\,4)(5\,6)$ negates the entire form.

Every endpoint is therefore a strict local maximum on $C_6$, by
Theorem~\ref{thm:sixdim-golden-local}. A nonconstant segment of global
maximizers cannot end at such a point. Working backward through the finite
merging path shows that the original form was the same endpoint. Thus every
maximizer is a seven-label relabeling of $\Omega_*$. Identical normalized
barycentric-gradient pairings determine simplices up to affine
symplectomorphism, proving the equality statement in
Theorem~\ref{thm:sixdim-global}.
